\chapter{Iterative methods}\label{ch:iterative}

\laheader{davidson.h, lanczos.h, conjgrad.h, gmres.h, diis.h, bisection.h}

\section{Purpose and main idea}\label{sec:iterative:purpose}

This chapter describes the header-only iterative algorithms of \LA{}:
\begin{itemize}
\item \code{davidson()} (\header{davidson.h}): a few lowest (or closest to a target) eigenpairs of a large real symmetric or
complex hermitian matrix, or right eigenpairs with real eigenvalues of a real non-symmetric matrix;
\item \code{lanczos()} (\header{lanczos.h}): a few lowest (or closest to a target) eigenpairs of a large hermitian matrix;
\item \code{conjgrad()} (\header{conjgrad.h}) and \code{gmres()} (\header{gmres.h}): solution of large linear systems $Ax=b$,
or of the normal equations $A\Tr Ax=A\Tr b$ for rectangular $A$ (least squares);
\item the class \code{DIIS} (\header{diis.h}): Pulay's DIIS convergence acceleration of a general fixed-point iteration;
\item \code{bisection_find()} and \code{interpolation_find()} (\header{bisection.h}): search in a sorted array.
\end{itemize}
All of them work for \code{T = float} and \code{std::complex<float>} as well (where the double versions support the
corresponding real or complex type); their default thresholds then come from \code{LA_precision<float>}
(section~\ref{sec:concepts:traits}).
Dense (LAPACK based) diagonalization, linear solvers and matrix functions are in chapter~\ref{ch:nonclass}; the small
subspace problems of \code{davidson} and \code{lanczos} are solved by the routines described there
(\code{diagonalize}, \code{gdiagonalize}, \code{diagonalize3}, \code{linear_solve}). The quasi-Newton GDIIS optimizer built on
the \code{DIIS} class is in chapter~\ref{ch:optimize}.

All functions are templates defined in the headers (all except \header{bisection.h} are included by \header{la.h}) and are
instantiated by the user's code. The matrix is a template parameter \code{Matrix} which is accessed only through a few member
functions, so that the matrix need not be stored at all: the matrix-vector product may be computed ``directly'' from other
data (direct CI in quantum chemistry, operators given by a formula, a product of several matrices, \dots). Any of the \LA{}
matrix classes providing these functions (\code{NRMat}, \code{NRSMat}, \code{SparseMat}, \code{SparseSMat}, \code{CSRMat},
\code{SemiSparseMat}, chapters~\ref{ch:mat}, \ref{ch:smat}, \ref{ch:sparse}) can also be used directly.

\section{Requirements on the matrix class}\label{sec:iterative:matrix}

The member functions used are listed in table~\ref{tab:iterative:matrix}; they must be \code{const}. Their semantics follow
the \LA{} matrix classes:
\begin{lareference}
\lafunc{int nrows() const;  int ncols() const}{Dimensions (any integer type). The eigensolvers require a square matrix
(\code{laerror} otherwise).}
\lafunc{bool issymmetric() const}{\code{true} if the matrix is real symmetric or complex hermitian. \code{davidson} chooses the
symmetric or non-symmetric algorithm by it; \code{lanczos} calls \code{laerror} if it is \code{false}. Note that
\code{NRMat::issymmetric()} always returns \code{false}, \code{NRSMat} and \code{SparseSMat} return \code{true}, \code{SparseMat}
returns its symmetric flag.}
\lafunc{const T* diagonalof(NRVec<T> &r, const bool divide=0, bool cache=false) const}{With \code{divide==false}: store the
diagonal $A_{ii}$ into \code{r} (allocated by the caller with \code{nrows()} elements) and return a pointer to the $n$ diagonal
elements, normally \code{&r[0]}. The eigensolvers call it with \code{cache=true}; a direct matrix may then compute the diagonal
once, keep it, and return a pointer to its own copy (the comment in davidson.h recommends this ``for efficiency in a direct
CI case''); the library classes ignore \code{cache}. With \code{divide==true}: divide \code{r} element-wise by the nonzero
diagonal elements (Jacobi preconditioning) and return \code{NULL}. For a rectangular matrix ($m\times n$, $m\ne n$) the
library classes use the diagonal of $A\Tr A$ (squared column norms, \code{r} with $n=$\code{ncols()} elements) instead, which is
the right (Jacobi) preconditioner for the normal equations used by \code{conjgrad} and \code{gmres} in that case; user matrix
classes should follow the same convention.}
\lafunc{void gemv(const T beta, NRVec<T> &r, const char trans, const T alpha, const NRVec<T> &x) const}{$r\leftarrow\beta r+
\alpha\,\mathrm{op}(A)\,x$ with $\mathrm{op}(A)=A$ for \code{trans=='n'} and $A\Tr$ for \code{'t'}; \code{r} is allocated by
the caller with the right size. With $\beta=0$ the previous contents of \code{r} must be ignored (it can be uninitialized).}
\end{lareference}

\begin{table}[htb]
\centering
\begin{tabular}{@{}l>{\raggedright\arraybackslash}p{2.3cm}>{\raggedright\arraybackslash}p{3.1cm}>{\raggedright\arraybackslash}p{2.9cm}@{}}
\toprule
 & \code{issymmetric} & \code{diagonalof} & \code{gemv} \\
\midrule
\code{davidson} & yes & \code{divide=false} & \code{'n'}, $\beta=0$ \\
\code{lanczos} & yes (must be true) & \code{divide=false}, only without \code{initguess} & \code{'n'}, $\beta=0$ and $\beta=1$ \\
\code{conjgrad} & -- & \code{divide=true}, if \code{precondition} & \code{'n'}; \code{'t'} if \code{doguess} or non-square \\
\code{gmres} & -- & \code{divide=true}, if \code{precondition} & \code{'n'}; \code{'t'} if \code{doguess} or non-square \\
\code{explicit_matrix} & -- & -- & \code{'n'}, $\beta=0$ \\
\bottomrule
\end{tabular}
\caption{Member functions of \code{Matrix} used by the iterative solvers; all of them also use \code{nrows()} and \code{ncols()}.}\label{tab:iterative:matrix}
\end{table}

\paragraph{Copy on demand.} The solvers keep the Krylov vectors as shallow copies of their work vectors
(sec.~\ref{sec:concepts:refcount}), so the vectors \code{r} passed to \code{gemv} and \code{diagonalof} frequently share their
data with stored vectors. A user implementation must therefore not write elements of \code{r} through \code{operator[]} or a raw
pointer before calling \code{r.copyonwrite()}; whole-object operations (\code{r=0.}, \code{r*=beta}, \code{r|=d}, \code{axpy},
and the \code{gemv} of the library classes) call it themselves. Forgetting it silently corrupts the Krylov space (in DEBUG
builds the non-const \code{operator[]} of a shared vector raises \code{laerror}). The input vector \code{x} is a const reference
and can be read freely. The library classes satisfy this.

\paragraph{Out-of-core storage.} \code{davidson}, \code{lanczos}, \code{gmres} and \code{DIIS} can keep their vectors in
scratch files instead of memory (parameter \code{incore}). They use \code{AuxStorage<T>} (\header{auxstorage.h},
chapter~\ref{ch:utilities}): a binary file with fixed-length records of raw elements, created by \code{mkostemp} as
\texttt{AUX\_XXXXXX} in \code{\$TMPDIR} (if set, otherwise in the current working directory) and unlinked at once, so no file
remains, even after a crash or an exception. Only one $n$-vector (two in
\code{davidson}) is then kept in memory per stored set besides the work vectors, at the price of reading all stored vectors in
every iteration. If a solver is left by an exception, the storage objects themselves are not freed (raw pointers), but the
disk space is released when the process ends.

\paragraph{Template argument deduction.} The element type \code{T} of \code{davidson} and \code{lanczos} is deduced from the
vector arguments. A literal \code{NULL} passed for one of the pointer parameters depending on \code{T} (\code{eivecs},
\code{initguess}, \code{target}) makes the deduction fail; either cast it (\code{(NRVec<double> *)NULL}, as in t.cc) or give
\code{T} explicitly: \code{davidson<double>(a,e,NULL,NULL,2)} (verified). Arguments left at their default values are not a
problem.

%-----------------------------------------------------------------------------------------------------------------------------
\section{Davidson diagonalization}\label{sec:iterative:davidson}

\laheader{davidson.h}

\begin{lareference}
\lafunc{template <typename T, typename Matrix> void davidson(const Matrix &bigmat, NRVec<T> &eivals, NRVec<T> *eivecs, const char *eivecsfile, int nroots=1, const bool verbose=0, const double eps=LA_precision<normtype>::eigensolver_eps(), const bool incore=1, int maxit=100, const int maxkrylov = 500, void (*initguess)(NRVec<T> &)=NULL, const typename LA_traits<T>::normtype *target=NULL)}{
Davidson's method \cite{iterative:davidson1975} for the \code{nroots} lowest eigenvalues of \code{bigmat}.
\begin{itemize}
\item \code{eivals}: on return the eigenvalues in its first \code{nroots} elements (ascending, or by increasing distance from
\code{*target}); it must have at least \code{nroots} elements (\code{laerror}), it is not resized.
\item \code{eivecs}: \code{NULL} or an array of at least \code{nroots} vectors receiving the normalized eigenvectors (by deep
copy, \code{|=}, so their previous size does not matter).
\item \code{eivecsfile}: \code{NULL} or a file name; the eigenvectors are then also written into this file as binary records
of $n$ elements of type \code{T}, record $i$ = root $i$ (an existing file is removed first), also if
\code{eivecs==NULL}.
\item \code{nroots}: number of roots; reduced silently to \code{maxkrylov-1} if larger.
\item \code{verbose}: print the eigenvalue estimates in each iteration and the convergence of each root to \code{std::cout}.
\item \code{eps}: convergence threshold (see below).
\item \code{incore}: keep the Krylov vectors $v_k$ and $Av_k$ in memory (\code{true}) or in two \code{AuxStorage} scratch files.
\item \code{maxit}: maximum number of iterations, i.e.\ of matrix-vector products after the initial one (restarts included).
\item \code{maxkrylov}: maximum dimension of the subspace; the subspace matrices are allocated with this dimension. When it
is reached, the method restarts (see below).
\item \code{initguess}: function filling its argument (allocated with $n$ elements, uninitialized) with the start vector,
which need not be normalized. By default the unit vector $e_j$ with $j$ the index of the lowest (real part of the) diagonal
element is used.
\item \code{target}: if not \code{NULL}, the subspace eigenpairs are sorted by $|\lambda-\mathtt{*target}|$ instead of
ascending order, i.e.\ the roots closest to the target value are sought.
\end{itemize}
Throws \code{laerror} if not converged within \code{maxit} iterations.}
\end{lareference}

\paragraph{Algorithm.} The method keeps the subspace vectors $v_k$ (not orthonormalized exactly, therefore the overlap matrix
$S_{jk}=v_j\cdot v_k$ is kept too), the products $Av_k$, and the subspace matrix $H_{jk}=v_j\cdot Av_k$. In each iteration the
generalized subspace eigenproblem $Hc=\lambda Sc$ is solved by \code{diagonalize} (symmetric case) or by \code{gdiagonalize}
with \code{sorttype=2} (real eigenvalues first, in ascending order; non-symmetric case, right eigenvectors; the imaginary parts
of complex eigenvalues are discarded). The roots are converged one after another: for the current root $\lambda$ with the
Ritz vector $y=\sum_kc_kv_k$ the residual $\rho=\sum_kc_k(Av_k-\lambda v_k)$ is formed and the correction vector is
$\delta_i=\rho_i/\mathrm{d}_i$ with $\mathrm{d}_i=A_{ii}-\lambda$, where $|\mathrm{d}_i|$ is limited from below by $0.1$
(keeping its sign). The correction is normalized and orthogonalized by Gram-Schmidt to all previous subspace vectors (the
projection on each $v_j$ is repeated while it removes more than 1\,\% of the norm) and appended to the subspace. Each iteration
costs one \code{gemv} with \code{'n'} and one \code{diagonalof}.

\paragraph{Convergence.} A root is considered converged when the coefficient of the most recently added subspace vector in its
subspace eigenvector, $|c_{\rm last}|$, is below \code{eps}; this is \emph{not} the residual norm, and the error of the
eigenvalue is typically much smaller than \code{eps}. After each iteration the method continues with the lowest root not yet
converged, and returns to a lower root if its $|c_{\rm last}|$ exceeds \code{eps} again. The method finishes successfully when
all \code{nroots} roots are converged, or (with a warning on \code{std::cout}) when the new correction vector vanishes before or
during the orthogonalization; in this case the current Ritz values are returned even if not all roots satisfy the criterion
(intended for tiny matrices). If $|c_{\rm last}|>100$ for the current root, ``Divergence in Davidson'' is printed and
\code{laerror} is called.

\paragraph{Restart.} When the subspace dimension reaches \code{maxkrylov}, the method restarts from the current approximation
of the lowest root alone, but only if no root has converged yet; otherwise it fails with \code{laerror}. For several roots
\code{maxkrylov} has therefore to be large enough to converge all of them without restart.

\paragraph{Element types.} \code{T=double} is the main case. For complex hermitian matrices (\code{T=std::complex<double>})
the method works, but according to the source comment ``does not converge so well for more than 1 root''; non-hermitian complex
matrices are not supported. The non-symmetric real case yields right eigenvectors belonging to real eigenvalues only.

The file has no header; it is read back with the \code{get} methods without dimensions, e.g.\ all roots as the rows of a matrix,
\code{NRMat<T> V(nroots,n); V.get(fd,false);}, or root by root with \code{NRVec<T>::get(fd,false)} (\code{fd} from
\code{open(eivecsfile,O_RDONLY)}). \code{AuxStorage(name)} cannot be used for reading, it recreates the file.


\begin{lareference}
\lafunc{template <typename T, typename Matrix> NRMat<T> explicit_matrix(const Matrix &bigmat)}{The explicit dense matrix of an
implicitly given one, built column by column by \code{gemv} on unit vectors (for debugging). \code{T} must be given:
\code{explicit_matrix<double>(a)}. Same as \code{reconstructmatrix} in chapter~\ref{ch:nonclass}.}
\end{lareference}

%-----------------------------------------------------------------------------------------------------------------------------
\section{Lanczos diagonalization}\label{sec:iterative:lanczos}

\laheader{lanczos.h}

\begin{lareference}
\lafunc{template <typename T, typename Matrix> void lanczos(const Matrix &bigmat, NRVec<T> &eivals, NRVec<T> *eivecs, const char *eivecsfile, int nroots=1, const bool verbose=0, const double eps=LA_precision<normtype>::eigensolver_eps(), const bool incore=1, int maxit=100, const int maxkrylov = 1000, void (*initguess)(NRVec<T> &)=NULL, const typename LA_traits<T>::normtype *target=NULL)}{
Lanczos method for the \code{nroots} lowest (or closest to \code{*target}) eigenvalues of a hermitian matrix. The parameters
have the same meaning as for \code{davidson}, with these differences: \code{bigmat.issymmetric()} must be \code{true}
(\code{laerror} otherwise); the start vector is normalized by the function; \code{incore=false} keeps only the Lanczos vectors
$v_k$ in one scratch file; the default \code{maxkrylov} is 1000.}
\end{lareference}

\paragraph{Algorithm.} Plain three-term Lanczos recurrence $w=Av_j-\beta_{j-1}v_{j-1}$, $\alpha_j=v_j\cdot w$,
$w\leftarrow w-\alpha_jv_j$, $\beta_j=|w|$, $v_{j+1}=w/\beta_j$, one \code{gemv} per iteration (with $\beta=1$ in the
\code{gemv} sense, accumulating into $-\beta_{j-1}v_{j-1}$). The vectors are not reorthogonalized, so that in long runs
copies of converged eigenvalues may appear (a known property of the method in finite precision). If $\beta_j=0$ (invariant
subspace), a random vector orthonormalized to all previous $v_k$ is used to continue. The tridiagonal matrix is diagonalized
in each iteration by \code{diagonalize3}. All $v_k$ are stored (in core or in the scratch file) to build the eigenvectors at the
end. If the diagonal elements $v_j\cdot Av_j$ have an imaginary part larger than $10^{-6}$, \code{laerror} reports a
non-hermitian matrix.

\paragraph{Convergence.} Converged when the changes of the \code{nroots} lowest (or closest to the target) eigenvalues
between two iterations are all $\le$ \code{eps}; this is checked once the dimension exceeds \code{nroots}. There is no restart:
the method fails with \code{laerror} when more than \code{maxit} iterations are needed or when the dimension reaches
\code{maxkrylov} (the effective limit is $\min(\mathtt{maxit},\mathtt{maxkrylov}-1)$ iterations).

%-----------------------------------------------------------------------------------------------------------------------------
\section{Conjugate gradients}\label{sec:iterative:conjgrad}

\laheader{conjgrad.h}

\begin{lareference}
\lafunc{template<typename T, typename Matrix> bool conjgrad(const Matrix &bigmat, const NRVec<T> &b, NRVec<T> &x, const bool doguess=true, const double tol=LA_precision<normtype>::conjgrad_tol(), const int itmax=1000, const bool verbose=true, bool issquare=1, const bool precondition=1)}{
Solves $Ax=b$ by the (Jacobi preconditioned) conjugate gradient method (the source refers to the book of Stoer and Bulirsch).
Returns \code{true} if converged, \code{false} after \code{itmax}+1 iterations without convergence; \code{x} then holds the
last iterate.
\begin{itemize}
\item \code{b}: right-hand side with \code{nrows()} elements; \code{x}: initial guess and solution, must have \code{ncols()}
elements (\code{laerror} otherwise).
\item \code{doguess}: overwrite \code{x} by the start vector $-A\Tr b$ (divided by the diagonal if \code{precondition}),
normalized to 1. With \code{false} the contents of \code{x} are the start vector.
\item \code{tol}: convergence threshold, see below.
\item \code{itmax}: maximum number of iterations; \code{verbose}: print the error in each iteration (default \code{true}).
\item \code{issquare}: \code{true} solves $Ax=b$ directly, which requires a symmetric positive definite matrix; \code{false}
applies CG to the normal equations $A\Tr Ax=A\Tr b$ (two \code{gemv} per iteration), usable for non-symmetric matrices.
It is set to \code{false} automatically for a rectangular matrix.
\item \code{precondition}: divide the residual by the diagonal (\code{diagonalof(r,true)}) in each iteration. The diagonal must
then be positive. For a square matrix with \code{issquare=false} the diagonal of $A$ (not of $A\Tr A$) is used, since
\code{diagonalof} cannot know that the normal equations are solved; this is not the proper preconditioner for $A\Tr A$ (the
Jacobi preconditioner would be the squared column norms) and may slow down or break the convergence if $A$ has non-positive
diagonal elements; use \code{precondition=false} then.
\end{itemize}}
\end{lareference}

\paragraph{Convergence.} The criterion is $|p|\le$ \code{tol} for the current search direction $p$ (an absolute criterion;
in the first iteration $p$ is the preconditioned residual). The function calls \code{laerror} (``singular matrix'') if
$p\cdot Ap$ (or $|Ap|^2$) or the residual product $r\cdot z$ vanishes. Only real \code{T} is supported (the scalars are
\code{double}).

%-----------------------------------------------------------------------------------------------------------------------------
\section{GMRES}\label{sec:iterative:gmres}

\laheader{gmres.h}

\begin{lareference}
\lafunc{template<typename T, typename Matrix> bool gmres(const Matrix &bigmat, const NRVec<T> &b, NRVec<T> &x, const bool doguess=1, const double eps=LA_precision<normtype>::gmres_eps(), const int MAXIT=50, const bool verbose=1, bool square=1, const bool precondition=1, int neustart=0, const int incore=1)}{
Solves $Ax=b$ by restarted GMRES with Givens rotations (after H.~R.~Schwarz, original implementation by M.~Warken). Returns
\code{true} if converged; otherwise \code{false}, with \code{x} holding the last iterate.
\begin{itemize}
\item \code{b}, \code{x}, \code{doguess}, \code{precondition}: as for \code{conjgrad}; the preconditioner $D^{-1}$ (Jacobi) is
applied from the left, i.e.\ GMRES minimizes $|D^{-1}(Ax-b)|$.
\item \code{eps}: relative threshold; converged when the estimated (preconditioned) residual norm is below
\code{eps}$\cdot\beta_0$, $\beta_0$ being the preconditioned residual norm of the start vector.
\item \code{MAXIT}: dimension of the Krylov space in one cycle; \code{MAXIT+1} vectors are stored.
\item \code{verbose}: print the residual estimate in each step (default \code{true}).
\item \code{square}: \code{false} solves the normal equations $A\Tr Ax=A\Tr b$; set automatically for a rectangular matrix.
\item \code{neustart}: maximum number of cycles (restarts); 0 means \code{nrows()/10}, and at least 10 cycles are always allowed.
\item \code{incore}: 0 stores the Krylov vectors in an \code{AuxStorage} scratch file.
\end{itemize}
On convergence the minimizer over the whole current Krylov space is returned. A breakdown ($|h_{k+1,k}|<0.01\,$\code{eps},
the Krylov space is invariant) ends the current cycle with the solution in that space; if it is not converged (singular
matrix), the next cycle starts from the new residual.
Calls \code{laerror} for a $1\times1$ system (\code{ncols()==1}), for incompatible sizes and for an exactly singular triangular
factor. The scalars are \code{double}, so only real \code{T} is supported.}
\lafunc{template<class T> void gmres_backsubstitute(const NRMat<T> &R, NRVec<T> &c, const NRVec<T> &d, const int k)}{Auxiliary:
solves the upper triangular system $R_{0..k,0..k}\,c=d$ (calls \code{c.copyonwrite()}; \code{c} must be allocated).}
\lafunc{template<typename T> void gmres_update(NRMat<T> &H, const NRVec<T> &cci, const NRVec<T> &ssi, NRVec<T> &d, NRVec<T> &c, const int k, NRVec<T> *v, AuxStorage<T> *st, const int incore, NRVec<T> &x)}{Auxiliary:
applies the Givens rotations $0..k$ to the Hessenberg matrix \code{H} (in place), solves $Rc=-d$ and adds $\sum_{i=0}^{k}c_iv_i$ to \code{x}.}
\end{lareference}

%-----------------------------------------------------------------------------------------------------------------------------
\section{DIIS convergence acceleration}\label{sec:iterative:diis}

\laheader{diis.h}

Direct inversion in the iterative subspace \cite{iterative:pulay1980,iterative:pulay1982}: given the last $m$ iterates $x_i$ of
some iterative procedure and their error vectors $e_i$, the extrapolated vector is $\sum_ic_ix_i$ with the coefficients
minimizing $|\sum_ic_ie_i|$ subject to $\sum_ic_i=1$, obtained from the linear system
\begin{equation}\label{eq:iterative:diis}
\begin{pmatrix}0&-1&\cdots&-1\\-1&B_{11}&\cdots&B_{1m}\\\vdots&\vdots&&\vdots\\-1&B_{m1}&\cdots&B_{mm}\end{pmatrix}
\begin{pmatrix}\lambda\\c_1\\\vdots\\c_m\end{pmatrix}=\begin{pmatrix}-1\\0\\\vdots\\0\end{pmatrix},\qquad B_{ij}=e_i\cdot e_j .
\end{equation}

\begin{lareference}
\lafunc{template<typename T, typename U> class DIIS}{\code{T} is the type of the iterated quantity, \code{U} the type of the
error vector, typically \code{NRVec}, \code{NRMat} or \code{NRSMat} (e.g.\ \code{T=NRSMat<double>} for a Fock or density
matrix and \code{U=NRMat<double>} for the commutator error). Both template arguments must be given (there is no default).
Requirements: \code{T} and \code{U} need \code{operator|=}, \code{clear()}, \code{axpy()}, \code{copyonwrite()} and the traits
\code{LA_traits<>::elementtype} and \code{normtype}; \code{U} also \code{dot()} and \code{norm()}. The element types of
\code{T} and \code{U} must be the same (the matrix of eq.~\eqref{eq:iterative:diis} is built over the element type of \code{U}
and converted to that of \code{T}). For out-of-core storage \code{T} and \code{U} must be \code{NRVec}, \code{NRMat} or
\code{NRSMat} of plain data (\code{AuxStorage}) and keep their size.}
\lafunc{DIIS()}{Uninitialized object (for arrays of \code{DIIS}); \code{setup()} must be called before use, otherwise
\code{extrapolate} calls \code{laerror}.}
\lafunc{DIIS(const int n, const bool core=1)}{Subspace of at most \code{n} stored iterates, kept in memory (\code{core=true}) or
in two \code{AuxStorage} scratch files.}
\lafunc{void setup(const int n, const bool core=1)}{(Re)initialize as by the constructor; the previous storage is released and
all stored vectors are forgotten. Can be called repeatedly.}
\lafunc{void reset()}{Forget all stored vectors (the next \code{extrapolate} starts a new subspace) but keep the allocated storage
and the dimension. Used to restart DIIS, e.g.\ after a step judged unreasonable (\code{GDIIS_Optimizer}, chapter~\ref{ch:optimize}).
After \code{reset()} the object behaves exactly as a freshly constructed one (tested in testoptimize.cc).}
\lafunc{typename LA_traits<U>::normtype extrapolate(T &vec, U &errvec, bool verbose=false, const Un diiseps=0, bool errvecout=false)}{
Stores deep copies of \code{vec} and \code{errvec} as the newest pair (when \code{n} pairs are already stored, the oldest is
dropped: circular buffer), updates $B$ with the products \code{errvec.dot(}$e_i$\code{)} and solves
eq.~\eqref{eq:iterative:diis} for the current number $m$ of stored pairs (\code{linear_solve} for \code{NRSMat},
chapter~\ref{ch:nonclass}). \code{vec} is overwritten by $\sum_ic_ix_i$, and with \code{errvecout=true} also \code{errvec} by
$\sum_ic_ie_i$ (the estimated error of the extrapolated vector). \code{diiseps} (type \code{LA_traits<U>::normtype}) is added
to the diagonal element $B_{mm}$ of the new vector as a regularization. \code{verbose} prints the solution vector
$(\lambda,c_1,\dots,c_m)$. Returns $|\mathtt{errvec}|$ of the input error vector. On the first call after construction, \code{setup} or \code{reset} the pair is only stored, \code{vec} and
\code{errvec} are unchanged, and the return value is 1.}
\end{lareference}

Typical use: in each iteration compute the new iterate $x$ and its error $e$ (e.g.\ $e=x_{\rm new}-x_{\rm old}$ for a
fixed-point iteration, or the commutator $FDS-SDF$ in SCF), call \code{diis.extrapolate(x,e)} and continue with the
extrapolated $x$. The stored copies are independent of the user's objects (deep copies by \code{|=}), so the user may modify
\code{vec} and \code{errvec} freely afterwards.

For complex \code{U} the matrix $B_{ij}=e_i^\dagger e_j$ is hermitian: \code{errvec.dot(}$e_i$\code{)}$=e_{\rm new}^\dagger e_i$ is
stored in the lower triangle of the complex \code{NRSMat}, which is hermitian in arithmetic (chapter~\ref{ch:smat}), and the
coefficients $c_i$ are complex (tested on a complex fixed-point iteration, in core and out of core).

%-----------------------------------------------------------------------------------------------------------------------------
\section{Search in sorted arrays}\label{sec:iterative:bisection}

\laheader{bisection.h}

The header is included by \header{la.h}; included alone it needs \code{laerror} declared before (it does not include
\header{laerror.h}).

\begin{lareference}
\lafunc{template<typename INDEX, typename COMPAR, typename SUBJECT> INDEX bisection_find(INDEX dm, INDEX hm, const SUBJECT *key, const SUBJECT *base, unsigned int lead_dimension_base, COMPAR (*cmp)(const SUBJECT *, const SUBJECT *))}{
Bisection search of \code{*key} among the elements \code{base[i*lead_dimension_base]}, $i=\mathtt{dm},\dots,\mathtt{hm}$
(inclusive), which must be sorted in ascending order with respect to \code{cmp}. \code{cmp(a,b)} returns 0 if equal, a positive
value if \code{*a} $>$ \code{*b}, negative otherwise; it is always called with the array element as the first argument.
\code{lead_dimension_base} allows searching e.g.\ the first column of a row-major matrix (\code{SUBJECT} is the element type, the
key is compared with one element only). Returns the index $i$ of a matching element, or \code{dm-1} if not found (also for
\code{dm>hm}). \code{INDEX} should be a signed type (for \code{dm=0} the failure value is $-1$).}
\lafunc{template<typename INDEX, typename DISTANCE, typename SUBJECT> int interpolation_find(INDEX dm0, INDEX high, const SUBJECT *x, const SUBJECT *base, unsigned int lead_dimension_base, DISTANCE (*dist)(const SUBJECT *, const SUBJECT *))}{
Interpolation search with the same layout and result convention (index or \code{dm0-1}). \code{dist(a,b)} is a signed distance
$a-b$ (positive if \code{*a} $>$ \code{*b}, zero if equal); the next probe is placed by linear interpolation, so the speed
depends on how well \code{dist} reflects the distribution of the keys. Calls \code{laerror} if an interpolated index falls
outside the current interval (inconsistent \code{dist}).}
\end{lareference}

%-----------------------------------------------------------------------------------------------------------------------------
\section{Examples}\label{sec:iterative:examples}

\subsection{Davidson with a user-defined direct matrix}

The matrix $H_{ij}=d_i\delta_{ij}+c\,(\delta_{i,j+1}+\delta_{i+1,j})$ is never stored; the class provides only the functions
of table~\ref{tab:iterative:matrix}. The program was compiled and run; the three eigenvalues agree with dense diagonalization
and the residual norms $|Hv-\lambda v|$ are $10^{-12}$--$10^{-10}$ for \code{eps=1e-8}.
\begin{lstlisting}
#include "la.h"
using namespace LA;

//H = diag(d) + c*(nearest neighbour coupling), never stored
class Chain {
	NRVec<double> d;
	double c;
public:
	Chain(int n, double cc) : d(n), c(cc) {for(int i=0; i<n; ++i) d[i]=i+0.1*i*i;}
	int nrows() const {return d.size();}
	int ncols() const {return d.size();}
	bool issymmetric() const {return true;}
	const double *diagonalof(NRVec<double> &r, const bool divide=false, bool cache=false) const
		{
		if(divide) {r.copyonwrite(); for(int i=0; i<r.size(); ++i) r[i]/=d[i]; return NULL;}
		r|=d; return &r[0]; //deep copy, r may share data with a stored Krylov vector
		}
	void gemv(const double beta, NRVec<double> &r, const char trans, const double alpha, const NRVec<double> &x) const
		{
		int n=d.size();
		if(beta==0) r=0.; else r*=beta;
		r.copyonwrite(); //before writing r[i]
		for(int i=0; i<n; ++i)
			{
			double t=d[i]*x[i];
			if(i>0) t+=c*x[i-1];
			if(i<n-1) t+=c*x[i+1];
			r[i]+=alpha*t;
			}
		}
};

int main()
{
const int n=2000, nroots=3;
Chain h(n,0.5);
NRVec<double> e(nroots);
NRVec<double> *v = new NRVec<double>[nroots];
davidson(h,e,v,NULL,nroots,false,1e-8);            //vectors in memory
std::cout << "eigenvalues "<<e;
davidson(h,e,v,NULL,nroots,false,1e-8,false);      //vectors in scratch files
lanczos(h,e,v,NULL,nroots,false,1e-10);            //the same by Lanczos
NRMat<double> hd = explicit_matrix<double>(h);     //for checking by dense diagonalization
delete[] v;
}
\end{lstlisting}
The \code{gemv} above ignores \code{trans} since the matrix is symmetric; for \code{conjgrad} or \code{gmres} on a
non-symmetric direct matrix, \code{'t'} has to be implemented as well (or \code{doguess=false} used with a square matrix).

\subsection{Linear systems and DIIS}

\begin{lstlisting}
NRMat<double> a(n,n); ...          //non-symmetric, diagonally dominant
NRVec<double> b(n), x(n); ...
bool ok=gmres(a,b,x,true,1e-10,50,false);           //x: solution
NRSMat<double> s(n); ...           //symmetric positive definite
ok=conjgrad(s,b,x,true,1e-10,1000,false);

//DIIS acceleration of a fixed-point iteration x=f(x) (from testoptimize.cc)
DIIS<NRVec<double>,NRVec<double> > diis(4);
NRVec<double> x(0.,n);
for(int it=0; it<maxit; ++it)
	{
	NRVec<double> xn=f(x);
	NRVec<double> err=xn-x;
	diis.extrapolate(xn,err);      //xn becomes the extrapolated vector
	x=xn;
	}
\end{lstlisting}

t.cc contains further calls: \code{davidson} on \code{NRMat}, \code{NRSMat}, \code{SparseMat} (symmetric and non-symmetric,
out of core), \code{lanczos} and \code{davidson} compared with dense diagonalization, and \code{gmres} on dense and sparse
matrices. testoptimize.cc tests the in-core and out-of-core \code{DIIS} and \code{reset()}.

%-----------------------------------------------------------------------------------------------------------------------------
\section{Pitfalls}\label{sec:iterative:pitfalls}

\begin{itemize}
\item A user \code{gemv} or \code{diagonalof} writing into \code{r} via \code{operator[]} without \code{r.copyonwrite()}
corrupts the stored Krylov vectors (sec.~\ref{sec:iterative:matrix}).
\item A literal \code{NULL} for \code{eivecs}, \code{initguess} or \code{target} breaks template argument deduction of
\code{davidson}/\code{lanczos}; cast it or write \code{davidson<double>(...)}.
\item \code{NRMat::issymmetric()} is \code{false}: \code{davidson} on an \code{NRMat} always uses the non-symmetric algorithm, and
\code{lanczos} refuses it. Use \code{NRSMat} (or a \code{SparseMat} with the symmetric flag) for symmetric matrices.
\item The Davidson threshold \code{eps} bounds the last expansion coefficient, the Lanczos one the change of the eigenvalues,
the CG one the norm of the search direction (absolute) and the GMRES one the residual relative to the initial one.
\item \code{davidson} with several roots fails if \code{maxkrylov} is reached after the first root has converged; \code{lanczos}
has no restart at all.
\item The diagonal used for preconditioning must have no small elements of the wrong sign (zero elements are skipped by the
library classes); in \code{davidson} the denominators $|A_{ii}-\lambda|$ are bounded by 0.1 from below.
\item \code{doguess=true} (the default of \code{conjgrad} and \code{gmres}) discards the contents of \code{x}.
\item Scratch files of the out-of-core variants are created in the current directory and survive an exception.
\item \code{DIIS} is not copyable (copy constructor and assignment are deleted, it owns raw storage); use references or pointers,
or default-construct it in arrays (\code{NRVec<DIIS<...> >}) and call \code{setup()}.
\end{itemize}
